\section{基于LQR的姿态最优控制}
\label{sec:lqr}

第\ref{sec:quat_control}节的四元数SAS采用手工整定的对角增益，第\ref{sec:INDI}节的INDI依赖在线角加速度估计，两者在tail-sitter大角度机动中的表现均存在局限（见第\ref{sec:vtol}节仿真）。本节在悬停工作点建立姿态误差动力学的线性模型，采用线性二次型调节器（LQR）系统化地设计全状态反馈增益，使控制效能矩阵$\mathbf{B}_{\mathrm{true}}$的跨通道耦合项被最优利用而非回避。

\subsection{悬停点姿态误差动力学线性化}
\label{sec:lqr_model}

取状态向量$\mathbf{x}=[\boldsymbol{\varepsilon}_{\mathrm{err}}^T,\ \boldsymbol{\omega}^T]^T\in\mathbb{R}^6$，输入向量$\mathbf{u}=[\Delta\omega,\ \delta_t,\ \delta_f]^T$。在小误差假设$\|\boldsymbol{\varepsilon}_{\mathrm{err}}\|\ll 1$下，误差四元数运动学线性化为
\begin{equation}
    \dot{\boldsymbol{\varepsilon}}_{\mathrm{err}} \approx \frac{1}{2}\left(\boldsymbol{\omega} - \boldsymbol{\omega}_{\mathrm{ref}}\right)
    \label{eq:lqr_eps_dot}
\end{equation}
其中$\boldsymbol{\omega}_{\mathrm{ref}}$为期望姿态轨迹的参考角速率（定点悬停时为零）。悬停点角速度近似为零，$\boldsymbol{\omega}\times\mathbf{I}\boldsymbol{\omega}$与陀螺耦合项为高阶小量，故角速率动力学线性化为
\begin{equation}
    \dot{\boldsymbol{\omega}} \approx \mathbf{I}^{-1}\,\mathbf{B}_{\mathrm{true}}\big|_{\delta=0}\ \mathbf{u}
    \label{eq:lqr_omega_dot}
\end{equation}
其中$\mathbf{I}=\operatorname{diag}(I_x,I_y,I_z)$，$\mathbf{B}_{\mathrm{true}}$为式(\ref{eq:B_true})在$\delta_f=\delta_t=\Delta\omega=0$处的取值。代入MODEL-DEFAULT参数（$\omega_0=761.4$\,rad/s）：
\begin{equation}
    \mathbf{B}_{\mathrm{true}}\big|_{\delta=0} =
    \begin{bmatrix}
        -2\tau_0 & 0 & 0 \\
        0 & -bT_0 & -\tau_0 \\
        0 & -\tau_0 & aT_0
    \end{bmatrix}
    \label{eq:B_hover}
\end{equation}
其中$T_0=k_T\omega_0^2$，$\tau_0=k_Q\omega_0^2$。注意俯仰行（第二行）含$-\tau_0$项（前摆的反扭矩耦合）、偏航行（第三行）含$-\tau_0$项（尾摆的反扭矩耦合）——手工对角整定丢弃了这些交叉项，而LQR会主动利用它们。式(\ref{eq:B_hover})行列式非零，$(\mathbf{A},\mathbf{B})$可控：
\begin{equation}
    \mathbf{A} = \begin{bmatrix} \mathbf{0}_3 & \frac{1}{2}\mathbf{I}_3 \\ \mathbf{0}_3 & \mathbf{0}_3 \end{bmatrix},
    \qquad
    \mathbf{B} = \begin{bmatrix} \mathbf{0}_3 \\ \mathbf{I}^{-1}\mathbf{B}_{\mathrm{true}} \end{bmatrix}
    \label{eq:lqr_AB}
\end{equation}

\subsection{LQR问题描述与权重设计}
\label{sec:lqr_design}

最小化无限时域二次型代价
\begin{equation}
    J = \int_0^{\infty}\left(\mathbf{x}^T\mathbf{Q}\mathbf{x} + \mathbf{u}^T\mathbf{R}\mathbf{u}\right)\mathrm{d}t
    \label{eq:lqr_cost}
\end{equation}
最优反馈$\mathbf{u}=-\mathbf{K}\mathbf{x}$的增益由连续代数Riccati方程（CARE）给出：
\begin{equation}
    \mathbf{A}^T\mathbf{S} + \mathbf{S}\mathbf{A} - \mathbf{S}\mathbf{B}\mathbf{R}^{-1}\mathbf{B}^T\mathbf{S} + \mathbf{Q} = \mathbf{0},
    \qquad
    \mathbf{K} = \mathbf{R}^{-1}\mathbf{B}^T\mathbf{S}
    \label{eq:care}
\end{equation}
权重按Bryson规则选取——以各状态/输入的最大允许值归一化：
\begin{equation}
    Q_{ii} = \frac{1}{x_{i,\max}^2},\qquad R_{ii} = \frac{1}{u_{i,\max}^2}
    \label{eq:bryson}
\end{equation}
取$\varepsilon_{\max}=0.5$（约$60^\circ$姿态误差）、$\omega_{\max}=1$\,rad/s、$\mathbf{u}_{\max}=[\Delta\omega_{\max},\ \delta_{\max},\ \delta_{\max}]$（来源: \texttt{models/aircraft-model.json}），得$\mathbf{Q}=\operatorname{diag}(4,4,4,1,1,1)$，$\mathbf{R}=\operatorname{diag}(2.04,5.25,5.25)$。

\subsection{增益矩阵与闭环极点分析}
\label{sec:lqr_poles}

数值求解CARE（仿真实现：\texttt{controllers.py}的\texttt{QuatLQRController}类）得
\begin{equation}
    \mathbf{K} =
    \begin{bmatrix}
        -1.400 & 0 & 0 & -0.747 & 0 & 0 \\
        0 & -0.867 & -0.102 & 0 & -0.474 & -0.056 \\
        0 & -0.102 & 0.867 & 0 & -0.056 & 0.476
    \end{bmatrix}
    \label{eq:lqr_K}
\end{equation}
非对角项$K_{23}=K_{32}=-0.102$即为对反扭矩交叉耦合的最优补偿——尾摆产生俯仰力矩时以前摆抑制其反扭矩副作用，反之亦然。闭环极点$\operatorname{eig}(\mathbf{A}-\mathbf{B}\mathbf{K})$为
\begin{equation}
    \left\{-14.39,\ -10.17,\ -9.60,\ -1.01,\ -1.01,\ -1.00\right\}\ \text{rad/s}
    \label{eq:lqr_poles}
\end{equation}
全部位于左半平面且均为实数：三个约$-1$\,rad/s的慢模态对应姿态角误差动态（由$\mathbf{Q}$中$\varepsilon$权重主导），三个$-9.6\sim-14.4$\,rad/s的快模态对应角速率动态。无复极点意味着闭环响应无振荡模态——这与第\ref{sec:vtol}节仿真中LQR无回摆的时域表现互为印证。

\subsection{大角度机动的参考前馈扩展}
\label{sec:lqr_ff}

LQR调节器（式\ref{eq:care}）针对定点镇定设计。slerp过渡等大角度机动中，直接将阻尼作用于绝对角速率$\boldsymbol{\omega}$会与期望旋转对抗，产生跟踪滞后。引入参考状态前馈，将反馈作用于跟踪误差：
\begin{equation}
    \mathbf{u} = -\mathbf{K}\begin{bmatrix} \boldsymbol{\varepsilon}_{\mathrm{err}} \\ \boldsymbol{\omega} - \boldsymbol{\omega}_{\mathrm{ref}} \end{bmatrix}
    \label{eq:lqr_tracking}
\end{equation}
其中$\boldsymbol{\omega}_{\mathrm{ref}}$由slerp轨迹解析给出（定轴旋转$\Rightarrow$机体系参考角速率为常值，见式\ref{eq:slerp}）。由于线性化模型式(\ref{eq:lqr_eps_dot})在跟踪误差小时沿整条过渡路径成立（$\mathbf{B}_{\mathrm{true}}$为机体系矩阵，与姿态无关），同一套增益无需调度即可覆盖从水平到竖直的全过程。

\subsection{仿真验证}
\label{sec:lqr_sim}

在3秒slerp垂直起飞工况下对比三种姿态控制架构（图\ref{fig:vtol_takeoff}、图\ref{fig:vtol_takeoff_compare}），定量指标见表\ref{tab:ctrl_compare}。LQR在slerp结束时刻（$t=3$\,s）即收敛至$\pm2^\circ$包络内，过渡后最大回摆仅$1.96^\circ$；对比之下，手工整定的四元数SAS（含阻尼增益调度）回摆达$22.7^\circ$且10秒内未收敛，INDI（v2，含速率前馈）回摆$6.1^\circ$、6.1秒收敛。悬停抗扰工况（$3^\circ$初始俯仰偏差，图\ref{fig:vtol_ctrl_compare}）中LQR同样最快：5秒末$\|\boldsymbol{\varepsilon}_{\mathrm{err}}\|$降至$1.9\times10^{-4}$，比INDI低一个数量级。

\begin{table}[H]
\centering
\caption{垂直起飞工况三种姿态控制器性能对比（3\,s slerp, MODEL-DEFAULT参数）}
\label{tab:ctrl_compare}
\small
\begin{tabular}{l c c c c}
\toprule
\textbf{控制器} & \textbf{稳定时间 ($\pm2^\circ$)} & \textbf{最大回摆} & \textbf{$\varepsilon$ 峰值} & \textbf{$\varepsilon$ (10\,s末)} \\
\midrule
四元数SAS（增益调度+前馈） & $>10$\,s & $22.7^\circ$ & 0.198 & $4.4\times10^{-2}$ \\
四元数INDI（v2） & 6.12\,s & $6.1^\circ$ & 0.066 & $1.5\times10^{-3}$ \\
四元数LQR & \textbf{2.94\,s} & $\mathbf{2.0^\circ}$ & \textbf{0.017} & $\mathbf{2.1\times10^{-5}}$ \\
\bottomrule
\end{tabular}
\end{table}

结果印证了系统设计的价值：SAS的对角手工增益回避了效能矩阵的交叉耦合，INDI受角加速度估计带宽限制，而LQR通过CARE一次性最优平衡了跟踪精度、阻尼与执行器代价，并主动利用交叉耦合项。需要强调的是，LQR增益依赖精确的$\mathbf{B}_{\mathrm{true}}$与惯量参数——在MODEL-DEFAULT参数未标定的概念阶段，其优势是结构性的；参数摄动下的鲁棒性对比（LQR vs INDI）是后续工作。

\subsection{交叉项补偿的工作点局限性}
\label{sec:lqr_coupling_limit}

第\ref{sec:lqr_model}节的线性化在悬停工作点（$\delta_f=\delta_t=\Delta\omega=0$）冻结了效能矩阵$\mathbf{B}_{\mathrm{true}}|_{\delta=0}$（式\ref{eq:B_hover}）。这意味着LQR对交叉项的补偿是\textbf{开环前馈式}的——增益$\mathbf{K}$在离线设计时固化，运行中不随实际摆角更新。当执行器偏转使真实交叉项$-\tau_f\cos\delta_f$偏离设计值$-\tau_0$时，补偿精度退化。这与INDI形成对照：INDI每一步在线重构$\mathbf{B}_{\mathrm{true}}(\mathbf{u}_k)$（式\ref{eq:B_true}），交叉项在任意工作点保持精确。

为量化这一退化，在24\,m/s巡航配平状态下施加前摆$\delta_f$阶跃（产生俯仰耦合$-\tau_f\sin\delta_f$），扫描阶跃幅度，对比三种架构的俯仰耦合抑制能力（陀螺噪声$0.002$\,rad/s，5次蒙特卡洛均值）。结果见表\ref{tab:lqr_coupling_deg}。

\begin{table}[H]
\centering
\caption{前摆$\delta_f$阶跃下的俯仰耦合峰值（°，5次蒙特卡洛均值）}
\label{tab:lqr_coupling_deg}
\small
\begin{tabular}{l c c c c c}
\toprule
\textbf{控制架构} & $\delta_f=5^\circ$ & $10^\circ$ & $15^\circ$ & $20^\circ$ & $25^\circ$ \\
\midrule
SAS（对角分配） & 1.23 & 1.63 & 1.95 & 2.18 & 2.36 \\
LQR（冻结$\mathbf{B}|_{\delta=0}$） & 0.98 & 1.33 & 1.63 & 1.86 & 2.02 \\
INDI（在线$\mathbf{B}_{\mathrm{true}}$） & \textbf{0.40} & \textbf{0.40} & \textbf{0.35} & \textbf{0.44} & \textbf{0.48} \\
\bottomrule
\end{tabular}
\end{table}

结果清晰揭示了三种架构在交叉项处理上的层次与边界：

\begin{enumerate}[label=(\arabic*)]
    \item \textbf{LQR的耦合抑制随摆角幅度单调退化}：从$\delta_f=5^\circ$的$0.98^\circ$升至$25^\circ$的$2.02^\circ$，逐渐逼近对角SAS（$2.36^\circ$）。根源是冻结的$\mathbf{B}|_{\delta=0}$在大摆角下与真实$\mathbf{B}_{\mathrm{true}}$的偏差扩大——满偏时$\cos 25^\circ\approx0.906$，交叉项被低估约$9\%$，且主对角项同步失真（见表\ref{tab:B_comparison}）。因此LQR的交叉项优势是\textbf{工作点局部}的。
    \item \textbf{INDI的耦合抑制对摆角幅度不敏感}：全程稳定在$0.4^\circ$附近，因为其在线Jacobian在每个控制周期精确反映当前摆角下的交叉项。这是INDI相对LQR在大机动工况下的\textbf{结构性优势}——它不依赖工作点假设。
    \item \textbf{三者梯度与表\ref{tab:ctrl_compare}互补}：表\ref{tab:ctrl_compare}显示LQR在slerp大角度机动的跟踪精度最优（受益于全状态反馈的最优增益平衡），而本表显示INDI在持续大摆角的交叉项补偿最优（受益于在线Jacobian）。两者优势的工况不同：LQR胜在瞬态跟踪品质，INDI胜在工作点变化的耦合鲁棒性。
\end{enumerate}

\textbf{结论}：LQR通过$\mathbf{B}|_{\delta=0}$在工作点附近有效利用交叉项，但其离线冻结的增益使其补偿精度随摆角偏离工作点而退化；INDI的在线$\mathbf{B}_{\mathrm{true}}$重构提供了与工作点无关的精确交叉项补偿。在需要持续大摆角机动的应用场景，INDI的耦合鲁棒性优于LQR；在以定点跟踪为主的场景，LQR的最优增益平衡仍具优势。这一互补性提示了"工作点附近LQR + 大机动INDI"的增益调度混合架构作为潜在的发展方向。

\subsection{级联架构：姿态外环与带交叉项补偿的角速度内环}
\label{sec:cascade}

第\ref{sec:lqr_coupling_limit}节的分析表明，无论是冻结$\mathbf{B}$的LQR还是在线$\mathbf{B}$的INDI，在\textbf{单级全状态反馈直接输出执行器指令}的结构下，交叉项补偿都难以兼顾最优性与工作点鲁棒性。本节提出一种级联控制架构，将姿态控制与角速度控制解耦，使交叉项补偿在内环被专门处理，从而同时获得两种优势。

\subsubsection{架构设计}

级联结构将控制问题按时间尺度分离为两个环：

\begin{enumerate}[label=(\arabic*)]
    \item \textbf{外环（姿态环，慢）}：纯比例（或比例-积分）控制，将姿态误差映射为角速度指令
    \begin{equation}
        \boldsymbol{\omega}_{\mathrm{ref}} = -\mathbf{K}_p\,\boldsymbol{\varepsilon}_{\mathrm{err}} - \mathbf{K}_i\int\boldsymbol{\varepsilon}_{\mathrm{err}}\,\mathrm{d}t + \boldsymbol{\omega}_{\mathrm{ref}}^{\mathrm{ff}}
        \label{eq:cascade_outer}
    \end{equation}
    外环不直接操作执行器，只生成角速度指令，因此结构极简（纯P/PI），无需考虑执行器动力学与交叉耦合；
    \item \textbf{内环（角速度环，快）}：将角速度误差映射为虚拟角加速度，再经在线效能矩阵分配为执行器增量
    \begin{align}
        \boldsymbol{\nu} &= \mathbf{K}_\omega\left(\boldsymbol{\omega}_{\mathrm{ref}} - \boldsymbol{\omega}\right)
        \label{eq:cascade_nu}\\
        \mathbf{u}_{k+1} &= \mathbf{u}_k + \mathbf{B}_{\mathrm{true}}^{-1}(\mathbf{u}_k)\,\boldsymbol{\nu}
        \label{eq:cascade_alloc}
    \end{align}
    内环通过在线重构$\mathbf{B}_{\mathrm{true}}(\mathbf{u}_k)$（式\ref{eq:B_true}）实现精确的交叉项补偿——前摆对俯仰的反扭矩耦合$-\tau_f\cos\delta_f$在分配时被显式计入。
\end{enumerate}

\textbf{设计要点}：与INDI不同，内环的角速度指令跟踪\textbf{不依赖角加速度估计}（$\boldsymbol{\nu}$由角速度误差经比例生成，$\mathbf{B}_{\mathrm{true}}^{-1}$仅作分配而非动态逆），因此规避了INDI的$\dot{\boldsymbol{\omega}}$噪声放大问题；与LQR不同，外环增益$\mathbf{K}_p$仅作用于姿态-角速度运动学（近似积分器），与效能矩阵解耦，因此在线$\mathbf{B}_{\mathrm{true}}$的引入不破坏外环设计。

\subsubsection{仿真验证}

在24\,m/s巡航配平、前摆$\delta_f$阶跃的耦合场景下（同表\ref{tab:lqr_coupling_deg}设置），对比级联架构与LQR、INDI的俯仰耦合抑制能力。结果见表\ref{tab:cascade_coupling}。

\begin{table}[H]
\centering
\caption{级联架构与LQR/INDI的俯仰耦合峰值对比（°，5次蒙特卡洛均值）}
\label{tab:cascade_coupling}
\small
\begin{tabular}{l c c}
\toprule
\textbf{控制架构} & $\delta_f=10^\circ$ & $\delta_f=25^\circ$ \\
\midrule
纯LQR（冻结$\mathbf{B}$，单级全状态反馈） & 1.33 & 2.02 \\
纯INDI（在线$\mathbf{B}$，单级） & 0.40 & 2.86 \\
级联P外环 + 在线$\mathbf{B}_{\mathrm{true}}$内环 & \textbf{0.02} & \textbf{0.04} \\
级联PI外环 + 在线$\mathbf{B}_{\mathrm{true}}$内环 & \textbf{0.02} & \textbf{0.04} \\
\bottomrule
\end{tabular}
\end{table}

结果揭示了级联架构的显著优势及其物理根源：

\begin{enumerate}[label=(\arabic*)]
    \item \textbf{耦合抑制提升约两个数量级}：级联架构将俯仰耦合峰值从LQR的$1.33$--$2.02^\circ$、INDI的$0.40$--$2.86^\circ$降至$0.02$--$0.04^\circ$。时域上，$\delta_f=25^\circ$阶跃后俯仰角仅有$0.047^\circ$的微小瞬态即恢复，而LQR持续偏离$2.1^\circ$。
    \item \textbf{优势源于结构解耦而非增益提高}：单级全状态反馈中，交叉项补偿与姿态-角速度耦合动力学相互干扰，补偿力矩被系统动态"淹没"；级联结构将角速度环带宽提高到姿态环之上（时间尺度分离），内环快速抑制耦合扰动，使其无法传播到姿态环。这与第\ref{sec:control}节SAS的级联思想一致，但内环用在线$\mathbf{B}_{\mathrm{true}}$替代了对角近似。
    \item \textbf{纯P外环存在稳态误差}：由于外环仅比例反馈，常值耦合扰动下姿态存在小稳态偏差（$0.041^\circ$）；引入积分项（级联PI）后稳态误差降至$0.032^\circ$。该稳态误差的量级远小于LQR/INDI的耦合峰值，但对高精度姿态保持场景仍需积分校正。
    \item \textbf{INDI在大摆角下的异常}：纯INDI在$\delta_f=25^\circ$时耦合峰值（$2.86^\circ$）反超$10^\circ$（$0.40^\circ$），源于其增量指令限幅（$|\Delta\mathbf{u}|\le0.2$）在大摆角扰动下饱和，在线$\mathbf{B}_{\mathrm{true}}$的补偿能力被限幅钳制。级联架构通过提高内环带宽并用完整力矩分配（非增量限幅），避免了这一饱和。
\end{enumerate}

\subsubsection{与三种"LQR在线B"路径的关系}

在探索LQR交叉项补偿的在线化时，曾评估三种直接路径（详见\texttt{lqr\_online\_b.py}）：

\begin{itemize}
    \item \textbf{增益调度LQR}：离线多工作点CARE+在线插值。实质是离散逼近，且调度维度随摆角组合增长；
    \item \textbf{SDRE（在线Riccati）}：每周期在线解CARE（单次仿真1500次求解）。计算量大，且稳定性证明弱于LQR；
    \item \textbf{LQR外环+在线B分配}：保留离线LQR增益，分配环节用在线$\mathbf{B}_{\mathrm{true}}$求逆。
\end{itemize}

仿真表明三者均\textbf{未能}显著改善耦合抑制（峰值$1.3$--$2.6^\circ$，与纯LQR相当）。根本原因是它们都保留了\textbf{单级全状态反馈}结构——无论B矩阵如何更新，交叉项补偿仍与姿态-角速度耦合动力学竞争。级联架构的成功恰恰在于\textbf{改变了结构}（时间尺度分离），而非仅更新B矩阵。这一对比说明：\textbf{交叉项补偿的有效性更多取决于控制架构的解耦程度，而非效能矩阵的在线化程度}。

\subsubsection{结论}

级联"姿态P/PI外环 + 在线$\mathbf{B}_{\mathrm{true}}$角速度内环"架构，通过时间尺度分离将交叉项补偿从系统动态中解放出来，在本构型上实现了约两个数量级的耦合抑制改善，且不引入INDI的角加速度估计噪声问题、不依赖LQR的工作点假设。这一架构将第\ref{sec:control}节SAS的级联思想与第\ref{sec:INDI}节的在线效能矩阵相结合，是概念阶段向工程实现过渡的推荐控制结构。其代价是内环需要在线$3\times3$矩阵求逆（与INDI相同的计算量），以及外环积分项的引入。
